\begin{pastexam}[label={px:ncku-112-ee-ds-07}]{成大電機所 112 學年度 第 7 題}
（問答）

(15\%) A word-chain automaton is a directed graph with a finite number of vertices V, satisfying the
following properties. (i) There is a specified initial node $v_0\in \boldsymbol{V}$. (ii) There is
subset \textbf{F} of set $V$ which contains terminal nodes. (iii) Each edge is labeled with a symbol
from a finite set \textbf{W}. A string \textbf{w$_1$w$_2$\ldots w$_k$} of symbols in \textbf{W} is
recognized by the automaton if there is a directed path starting at the initial node $v_0$ and
ending at a terminal node w $\in \boldsymbol{F}$ such that the sequence of edges in this directed
path is \textbf{w$_1$w$_2$\ldots w$_k$}. The set \textbf{L} of all strings recognized by the
automaton is called the language recognized by the automaton. A word-chain automaton is a special
type of finite-state automaton, which in turn, is a type of finite-sate machine.

(a) Draw a word-chain automaton such that the following five sentences can be recognized: 1. I like
his rabbit. 2. They hate his cat. 3. I think that his hamster is cute. 4. I know that his rabbit
is soft. 5. They hate to have a hamster.

(b) Explain how the above automaton is used to check whether a sentence is recognized or not.
\end{pastexam}

\begin{pastsoln}
本題和成大智慧資安學程同年的第~7~題相同（電機所、電通所、智慧資安學程三個系所共用同一份考卷），完整的詳解見歷屆試題~\ref{px:ncku-112-isec-ds-07}，這裡只列答案和重點。

答案：
\begin{itemize}
\item (a) 把單字當成符號：從起點 $v_0$ 依句子的單字一路連邊，共用相同的開頭（I、They 開頭各一個分支；
I 之後分出 like、think、know），五個句子的結尾都接到同一個終點（雙圈）。 這樣恰好辨識這五個句子，
圖見智慧資安學程的那一題。
\item (b) 從 $v_0$ 出發，依序讀單字，每讀一個就沿著標籤是這個單字的邊走一步；找不到這樣的邊就拒絕；
讀完時停在終點集合 $F$ 裡就接受，否則拒絕。 每個節點的出邊用雜湊表存時，長度 $k$ 的句子檢查時間
是期望 $O(k)$，和資結書第~27.4~節用有限自動機比對字串是同一個做法。
\end{itemize}
\end{pastsoln}
